\Needspace{4\baselineskip}
\begin{Example}
\label{ex:abel-integral-equation-laplace}
Abel\textquotesingle s integral equation is
\begin{equation}\label{eq:abel-integral-equation}
\int_0^x\dfrac{y(t)}{\sqrt{x-t}}\,\mathrm{d}t=\sqrt\pi f(x),
\end{equation}
again we seek a solution for $x>0$ and we think of $y$ and $f$ as being extended by zero for negative $x$. In order to have a solution, we must obviously have $f(0)=0$. The left-hand side is the convolution of $y$ and $x_+^{-1/2}$, and the Laplace transform of a convolution is the product of the Laplace transforms. To find the transform of $x_+^{-1/2}$, we compute
\begin{equation}\label{eq:laplace-inverse-square-root-kernel}
\int_0^\infty e^{-sx}x^{-1/2}\,\mathrm{d}x=\dfrac1{\sqrt s}\int_0^\infty e^{-t}t^{-1/2}\,\mathrm{d}t=\sqrt{\dfrac\pi s}
\end{equation}
for any real positive $s$ and because of the uniqueness of analytic continuation this also holds for complex $s$. Hence the transformed equation reads
\begin{equation}\label{eq:abel-equation-laplace-transform}
\dfrac{\mathcal{L}[y](s)}{\sqrt s}=\mathcal{L}[f](s),
\end{equation}
which we write as
\begin{equation}\label{eq:abel-equation-transform-solution}
\mathcal{L}[y](s)=\dfrac{s\mathcal{L}[f](s)}{\sqrt s}.
\end{equation}
Transforming back, we find
\begin{equation}\label{eq:abel-equation-inverse-transform-solution}
y(x)=\dfrac1{\sqrt\pi}\int_0^x\dfrac{f'(t)}{\sqrt{x-t}}\,\mathrm{d}t.
\end{equation}

\end{Example}

\Needspace{4\baselineskip}
\begin{Example}
\label{ex:heat-boundary-laplace}
The Laplace transform is also applicable to initial-value problems for PDEs. We first remark that Definition~\ref{def:tempered-fourier-transform} is easily generalized to define the Fourier transform of a generalized function with respect to only a subset of the variables. For example, when dealing with an initial-value problem, we can take the Laplace transform with respect to time. Of course, to make sense of boundary conditions, one needs to know more about the solution than that it is a generalized function. For example, in the following problem, we may think of $u$ as a generalized function of $t$ depending on $x$ as a parameter.

We consider the initial/boundary-value problem
\begin{equation}\label{eq:heat-unit-boundary-zero-initial-problem}
\begin{aligned}
u_t&=u_{xx},\quad x\in(0,1),\ t>0,\\
u(x,0)&=0,\quad x\in(0,1),\\
u(0,t)&=u(1,t)=1,\quad t>0.
\end{aligned}
\end{equation}
As usual, we extend $u$ by zero for negative $t$. Laplace transform in time leads to the problem
\begin{equation}\label{eq:heat-unit-boundary-laplace-problem}
\begin{aligned}
s\mathcal{L}[u](x,s)&=\mathcal{L}[u]_{xx}(x,s),\\
\mathcal{L}[u](0,s)&=\mathcal{L}[u](1,s)=\dfrac1s.
\end{aligned}
\end{equation}
This equation has the solution
\begin{equation}\label{eq:heat-unit-boundary-laplace-solution}
\mathcal{L}[u](x,s)=\dfrac{\cosh(\sqrt s(x-1/2))}{s\cosh(\sqrt s/2)}.
\end{equation}
The formula for the inverse transform yields
\begin{equation}\label{eq:heat-unit-boundary-inverse-laplace-formula}
u(x,t)=\dfrac1{2\pi i}\int_{\mu-i\infty}^{\mu+i\infty}e^{st}\dfrac{\cosh(\sqrt s(x-1/2))}{s\cosh(\sqrt s/2)}\,\mathrm{d}s.
\end{equation}
Here we can take $\mu$ to be any positive number. The integral cannot be evaluated in closed form, but of course it can be evaluated numerically; it can also be used to deduce qualitative properties of the solution such as its regularity or its asymptotic behavior as $t\to\infty$.
\end{Example}
